\section{Related Work}

\subsection{Traditional Light Field Depth Estimation Methods}

Many traditional depth estimation methods for light field (LF) cameras have been proposed, primarily leveraging physical cues such as Epipolar Plane Image (EPI) slopes, defocus blur, and stereo matching [14]. These approaches substantially rely on geometric modeling under the strict Lambertian reflectance assumption, where a surface point emits consistent radiance across all viewpoints, forming distinct linear structures in the EPI space.

Wanner et al. [15] pioneered the use of EPI structural analysis by deriving a theory that links the local orientation of EPI lines to scene depth. Specifically, they employed a structure tensor to calculate the dominant slope of EPI lines, effectively resolving the ambiguity in textured regions. While this approach achieves notable accuracy on Lambertian surfaces, it intrinsically assumes that EPI lines are perfectly straight. Consequently, when encountering non-Lambertian surfaces such as specular reflections, the EPI lines distort into X-shaped patterns, causing the structure tensor-based slope calculation to fail completely.

To mitigate the limitations of purely slope-based methods, Tao et al. [16] extended the physical modeling by combining defocus and correspondence cues. They formulated a depth estimation framework that leverages the depth-dependent blur in focal stacks and the angular shift in sub-aperture images. By fusing these two complementary cues, their method effectively improves depth accuracy in weakly textured regions. However, this fusion strategy heavily relies on the consistency between defocus and correspondence cues. In the presence of glossy or mirror-like surfaces, the specular highlights shift inconsistently across viewpoints, leading to severe conflicts between the two cues and causing the algorithm to hallucinate their depths.

In contrast to focal stack-based approaches, stereo matching paradigms focus on leveraging multi-view geometry directly from sub-aperture images. Jeon et al. [17] proposed an adaptive cost volume method that incorporates a phase-shift technique to compensate for sub-pixel displacements. This approach effectively refines depth boundaries and resolves the ambiguity in fine structures. Nonetheless, it primarily depends on the photo-consistency assumption across viewpoints. When applied to spatially-varying (SV) reflectance or transparent or mirror (ToM) surfaces, the matching cost becomes under-constrained, leading to severe mismatches in non-Lambertian regions.

Despite their extensive contributions, these traditional approaches share inherent deficiencies that restrict their applicability in real-world scenarios. On the one hand, they heavily rely on hand-crafted features and the strict Lambertian assumption. When confronting non-Lambertian materials (e.g., specular reflections and highlights), the EPI geometry degenerates from straight lines to X-shaped patterns, rendering the slope and matching calculations entirely invalid. On the other hand, they lack the physical-level decoupling capability among illumination, depth, and material reflectance. As a result, they exhibit significantly poor robustness in complex textured and mixed-reflection scenes, failing to provide reliable priors for modern complex scenarios. These inherent limitations not only provide the motivation for our subsequent quantitative analysis of the physical mechanisms behind EPI failure but also underscore the necessity of explicitly decoupling geometric and material properties. To this end, our proposed unified dual-mask physical model addresses these gaps by introducing a three-layer angular signal decomposition and a geometric dual-mask routing mechanism, enabling adaptive processing of complex reflectance characteristics without being constrained by the traditional EPI assumptions.

\subsection{Deep Learning-based General Epipolar Depth Estimation Methods}

Following the limitations of handcrafted physical cues, deep learning-based methods have emerged to implicitly learn the mapping from Epipolar Plane Images (EPIs) to depth maps. Shin et al. [18] pioneered this direction by proposing EPINET, a fully convolutional network that extracts both spatial and angular features from EPIs and the central view. By concatenating these multi-dimensional features, the network effectively regresses depth values, achieving considerable accuracy on standard Lambertian benchmarks. Subsequent improvements [19] introduced multi-scale EPI processing to enhance local geometric representations. However, the reliance on standard convolutional kernels inherently restricts the receptive field along the epipolar lines, making it challenging to decouple and aggregate long-range geometric dependencies in high-resolution light fields.

To overcome the restricted receptive field and mitigate feature interference, recent studies have introduced attention mechanisms and feature disentanglement strategies. For instance, Wang et al. proposed a disentangled feature network that explicitly separates angular and spatial representations. By employing spatially-varying (SV) attention modules, this approach adaptively aggregates multi-view cues, significantly enhancing depth estimation in texture-less and occluded regions. Similarly, Meng et al. [20] leveraged an adaptive cost volume to regularize the distribution of epipolar features, further improving robustness against local ambiguities. While these attention-driven architectures successfully expand the effective receptive field and refine feature aggregation, they primarily focus on optimizing the network capacity without reconsidering the underlying physical assumptions of the EPI structure.

Further pushing the boundaries of receptive field expansion, Transformer-based architectures have been explored to capture global epipolar dependencies. Liang et al. [21] introduced a light field Transformer that applies self-attention mechanisms specifically along the epipolar lines. This design allows the model to dynamically weigh distant pixels within the same EPI, effectively resolving the aliasing bimodal distribution and slope ambiguities in complex textures. Consequently, empirical evaluations on public datasets demonstrate that such Transformer-based models achieve state-of-the-art performance. Nonetheless, the dramatic increase in computational complexity and the implicit reliance on the linear EPI prior remain critical bottlenecks when generalizing to real-world scenarios with diverse material properties.

Despite these extensive architectural advancements, general deep learning-based EPI-based methods share inherent physical limitations. On the one hand, the vast majority of these networks implicitly assume that EPIs possess perfect linear structures. When the input contains non-Lambertian pixels, the networks cannot distinguish whether the EPI slope blur is caused by continuous depth variations or material reflections, leading to a dramatic performance degradation in non-Lambertian regions. On the other hand, these methods lack explicit physical constraints and material awareness. They attempt to fit the data merely by increasing network complexity (e.g., deeper networks or attention mechanisms), which cannot overcome the physical performance boundary caused by the violation of the EPI assumption. This proves that pure parameter tuning cannot resolve physical assumption conflicts. In contrast, our proposed unified dual-mask physical model explicitly addresses these physical conflicts by formulating a three-layer angular signal decomposition and a geometric dual-mask routing mechanism. This design inherently decouples geometric and material properties, enabling adaptive processing of complex reflectance characteristics without being constrained by the implicit linear EPI assumptions of existing deep networks.

\subsection{Summary}

In summary, both traditional and deep learning-based light field depth estimation methods exhibit inherent limitations when handling non-Lambertian surfaces. Traditional methods heavily depend on hand-crafted features and strict Lambertian assumptions, failing completely when EPI geometries degenerate into X-shaped patterns. Conversely, deep learning-based approaches attempt to overcome these issues through increased network complexity, attention mechanisms, and global receptive fields. However, they still implicitly rely on linear EPI priors and lack explicit physical constraints, leading to suboptimal generalization and significant performance degradation in mixed-reflection scenes. Crucially, neither paradigm effectively decouples illumination, depth, and material reflectance at the physical level. To bridge this critical gap, this paper proposes a unified dual-mask physical model that explicitly formulates the physical mechanisms of EPI degradation and adaptively routes angular signals, thereby achieving robust and accurate depth estimation in complex real-world scenarios. Building upon this conceptual foundation, we now detail the complete technical pipeline, mathematical formulations, and implementation specifics of the proposed framework.